\chapter{Two Pointers and Subarrays}
\chapterepigraph{The sliding window's cousin, the two-pointer method, and
its algebraic sibling, the prefix sum, solve the ``count or find a
subarray with a property'' family. When the property grows monotonically
as the window widens, two moving pointers suffice; when it does not, a
prefix sum plus a hash map turns the search into counting equal keys.
These problems from the Sorting and Searching section round out the
window toolkit.}

\section{Two Pointers vs Prefix Sums}
\label{sec:cses-tp-intro}

Both techniques avoid examining all $O(n^2)$ subarrays, but they apply in
different regimes.

\begin{patternbox}[Which Technique for Which Subarray Problem?]
\begin{itemize}
  \item \textbf{Two pointers} -- when a window's property is
        \emph{monotonic}: widening the right end only ever increases the
        sum (positive values) or the distinct count, so the left end can
        chase it without ever moving backward. Used for Subarray Sums I,
        Subarray Distinct Values, Playlist.
  \item \textbf{Prefix sums + hash map} -- when values can be negative or
        the target is exact equality, monotonicity is lost; instead count
        how many earlier prefixes give the needed difference or residue.
        Used for Subarray Sums II, Subarray Divisibility.
  \item \textbf{Running best (Kadane)} -- when maximising a subarray sum,
        a single scan tracking ``best ending here'' suffices: Maximum
        Subarray Sum.
  \end{itemize}
\end{patternbox}

\begin{ideabox}[Prefix Sums Turn Subarray Questions into Pair Questions]
The sum of $a[l..r]$ is $\textit{pre}[r]-\textit{pre}[l-1]$. So ``a subarray
with sum $x$'' becomes ``two prefixes differing by $x$,'' countable in one
pass with a hash map of prefix values. This transform is the backbone of
the exact-sum and divisibility problems.
\end{ideabox}

\newpage

\section{Subarray Sums I}
\label{sec:cses-subarray-sums-i}

\problemstatement{Given an array of $n$ \emph{positive} integers and a
target $x$, count the number of contiguous subarrays whose sum equals $x$.}

\begin{patternbox}
With all values positive, the window sum is \textbf{monotonic}: extending
the right end only increases it, and advancing the left end only decreases
it. That monotonicity is exactly what lets a \textbf{two-pointer} sliding
window count matches in one linear pass.
\end{patternbox}

\subsection*{Why two pointers work here}

Keep a window $[l,r]$ and its sum. Advance $r$, adding $a[r]$. Whenever the
sum exceeds $x$, advance $l$, subtracting $a[l]$ -- safe because positivity
guarantees shrinking only lowers the sum. Whenever the sum equals $x$,
record one subarray. Because the sum never decreases when $r$ moves right
and never increases when $l$ moves right, each pointer only moves forward,
giving $O(n)$.

\begin{ideabox}[Positivity Makes the Window Monotone]
Two pointers require that growing the window changes the tracked quantity
in one direction only. Positive values guarantee that for sums -- which is
why Subarray Sums~II (negatives allowed) must abandon two pointers for
prefix sums.
\end{ideabox}

\subsection*{Algorithm}

\begin{enumerate}
  \item $l=0$, $\textit{sum}=0$, $\textit{count}=0$.
  \item For $r=0\ldots n-1$: add $a[r]$; while $\textit{sum}>x$ advance
        $l$ subtracting $a[l]$; if $\textit{sum}=x$ increment
        $\textit{count}$.
  \item Output $\textit{count}$.
\end{enumerate}

\subsection*{Dry run}

\dryrun{$a=[2,4,1,2,7]$, $x=7$.}

\begin{itemize}
  \item $[2,4,1]$ sums to $7\Rightarrow$ count $1$.
  \item Window shrinks/extends; $[4,1,2]=7\Rightarrow$ count $2$; later
        $[7]=7\Rightarrow$ count $3$.
  \item Answer $3$.
\end{itemize}

\subsection*{C++ implementation}

\begin{lstlisting}
#include <bits/stdc++.h>
using namespace std;

int main() {
    int n; long long x;
    cin >> n >> x;
    vector<long long> a(n);
    for (auto& v : a) cin >> v;

    long long sum = 0, count = 0;
    int l = 0;
    for (int r = 0; r < n; r++) {
        sum += a[r];
        while (sum > x) sum -= a[l++];         // positivity: safe to shrink
        if (sum == x) count++;
    }
    cout << count << "\n";
    return 0;
}
\end{lstlisting}

\begin{complexitybox}
\textbf{Time Complexity.} $O(n)$ -- both pointers only move forward.
\textbf{Space Complexity.} $O(1)$.
\end{complexitybox}

\begin{takeawaybox}
Two pointers count target-sum subarrays in linear time precisely because
positive values keep the window sum monotone. The moment negatives enter,
monotonicity breaks and the prefix-sum-plus-hash-map method takes over.
\end{takeawaybox}

\section{Subarray Sums II}
\label{sec:cses-subarray-sums-ii}

\problemstatement{Given an array of $n$ integers (which may be negative)
and a target $x$, count the number of contiguous subarrays whose sum equals
$x$.}

\begin{patternbox}
With negatives, the window sum is no longer monotone, so two pointers fail.
Instead use \textbf{prefix sums plus a hash map}: a subarray $[l+1,r]$ has
sum $x$ iff $\textit{pre}[r]-\textit{pre}[l]=x$, so count, for each prefix,
how many earlier prefixes equal $\textit{pre}[r]-x$.
\end{patternbox}

\subsection*{Counting equal prefix differences}

Let $\textit{pre}[i]$ be the sum of the first $i$ elements. A subarray
ending at position $r$ with sum $x$ corresponds to an earlier prefix
$\textit{pre}[l]=\textit{pre}[r]-x$. Sweep left to right maintaining a hash
map from prefix-sum value to how many times it has occurred; at each step
add the map's count of $\textit{pre}[r]-x$ to the answer, then record
$\textit{pre}[r]$. Initialise the map with $\{0:1\}$ so subarrays starting
at index $0$ are counted.

\begin{ideabox}[Turn a Subarray Sum into a Prefix Difference]
Any subarray sum is a difference of two prefix sums, so ``sum equals $x$''
becomes ``two prefixes differ by $x$.'' A hash map of previously seen prefix
values counts those pairs in one pass, negatives and all.
\end{ideabox}

\subsection*{Algorithm}

\begin{enumerate}
  \item Map $\{0:1\}$, $\textit{pre}=0$, $\textit{count}=0$.
  \item For each element: $\textit{pre}\mathrel{+}=a[i]$; add
        $\textit{map}[\textit{pre}-x]$ to $\textit{count}$; increment
        $\textit{map}[\textit{pre}]$.
  \item Output $\textit{count}$.
\end{enumerate}

\subsection*{Dry run}

\dryrun{$a=[2,-1,3,1]$, $x=4$.}

\begin{itemize}
  \item Prefixes: $0,2,1,4,5$. Looking for earlier prefix
        $\textit{pre}-4$.
  \item At $\textit{pre}=4$: $\textit{pre}-4=0$ seen once $\Rightarrow$
        count $1$ ($[2,-1,3]$). At $\textit{pre}=5$: $5-4=1$ seen once
        $\Rightarrow$ count $2$ ($[3,1]$).
  \item Answer $2$.
\end{itemize}

\subsection*{C++ implementation}

\begin{lstlisting}
#include <bits/stdc++.h>
using namespace std;

int main() {
    int n; long long x;
    cin >> n >> x;
    unordered_map<long long,long long> seen;
    seen[0] = 1;                               // empty prefix

    long long pre = 0, count = 0;
    for (int i = 0; i < n; i++) {
        long long a; cin >> a;
        pre += a;
        auto it = seen.find(pre - x);          // earlier prefix that closes a sum-x subarray
        if (it != seen.end()) count += it->second;
        seen[pre]++;
    }
    cout << count << "\n";
    return 0;
}
\end{lstlisting}

\begin{complexitybox}
\textbf{Time Complexity.} $O(n)$ amortised with a hash map. \textbf{Space
Complexity.} $O(n)$ for the map.
\end{complexitybox}

\begin{takeawaybox}
Once negatives break monotonicity, prefix sums plus a hash map count
exact-sum subarrays in one pass -- the go-to when two pointers cannot apply.
This is the same machinery behind ``subarray with sum divisible by $n$''
(next section), swapping the key from a value to a residue.
\end{takeawaybox}

\section{Subarray Divisibility}
\label{sec:cses-subarray-divisibility}

\problemstatement{Given an array of $n$ integers, count the number of
contiguous subarrays whose sum is divisible by $n$.}

\begin{patternbox}
``Sum divisible by $n$'' is prefix sums modulo $n$: a subarray sum is a
multiple of $n$ iff its two bounding prefix sums have the \textbf{same
residue} mod $n$. So count pairs of equal residues -- a hash map (or array)
of residue frequencies does it in one pass.
\end{patternbox}

\subsection*{Equal residues bound divisible subarrays}

$\textit{sum}(l+1,r)=\textit{pre}[r]-\textit{pre}[l]$ is divisible by $n$
exactly when $\textit{pre}[r]\equiv\textit{pre}[l]\pmod n$. So for each
residue class, any two prefixes in it form a valid subarray. Sweeping the
prefixes and tallying how many share each residue, the number of divisible
subarrays is $\sum_r\binom{c_r}{2}$ where $c_r$ counts prefixes with
residue $r$ (including the empty prefix, residue $0$). Handle negative sums
by normalising the residue into $[0,n)$.

\begin{ideabox}[Residue Collisions Are Divisible Subarrays]
Reducing prefix sums mod $n$ converts ``divisible difference'' into ``equal
residue.'' Counting how many prefixes fall in each residue class, then
choosing $2$ from each class, counts every divisible subarray without
examining pairs directly.
\end{ideabox}

\subsection*{Algorithm}

\begin{enumerate}
  \item Frequency array $\textit{cnt}[0\ldots n-1]$ with
        $\textit{cnt}[0]=1$ (empty prefix).
  \item Running prefix mod $n$ (normalised non-negative); for each element
        add $\textit{cnt}[\textit{residue}]$ to the answer, then increment
        that residue's count.
  \item Output the accumulated count.
\end{enumerate}

\subsection*{Dry run}

\dryrun{$a=[1,2,3,4,5]$, $n=5$.}

\begin{itemize}
  \item Prefixes mod $5$: $0,1,3,1,0,0$.
  \item Residue $0$ occurs at prefixes (empty),$4$,$5$ -- three of them;
        residue $1$ twice.
  \item Answer $\binom{3}{2}+\binom{2}{2}=3+... $ counted incrementally to
        the correct total of divisible subarrays.
\end{itemize}

\subsection*{C++ implementation}

\begin{lstlisting}
#include <bits/stdc++.h>
using namespace std;

int main() {
    int n;
    cin >> n;
    vector<long long> cnt(n, 0);
    cnt[0] = 1;                                // empty prefix has residue 0

    long long pre = 0, ans = 0;
    for (int i = 0; i < n; i++) {
        long long a; cin >> a;
        pre = ((pre + a) % n + n) % n;         // normalise residue to [0,n)
        ans += cnt[pre];                       // earlier prefixes with same residue
        cnt[pre]++;
    }
    cout << ans << "\n";
    return 0;
}
\end{lstlisting}

\begin{complexitybox}
\textbf{Time Complexity.} $O(n)$. \textbf{Space Complexity.} $O(n)$ for the
residue counts.
\end{complexitybox}

\begin{takeawaybox}
Divisibility of a subarray sum is equality of prefix residues; count
residue collisions in one pass. This is Subarray Sums~II with the map key
changed from a raw prefix value to its residue mod $n$ -- a small twist on
the same prefix-and-count template.
\end{takeawaybox}

\section{Subarray Distinct Values}
\label{sec:cses-subarray-distinct}

\problemstatement{Given an array of $n$ integers and an integer $k$, count
the number of contiguous subarrays that contain at most $k$ distinct
values.}

\begin{patternbox}
``At most $k$ distinct'' is monotone -- shrinking a window never increases
its distinct count -- so a \textbf{variable-size two-pointer window} works:
for each right end, advance the left end until at most $k$ distinct remain,
and every subarray ending here within the window is valid.
\end{patternbox}

\subsection*{Count subarrays ending at each right end}

Maintain a window $[l,r]$ with a frequency map and its distinct count.
Extend $r$; if adding $a[r]$ pushes the distinct count above $k$, advance
$l$ (removing elements) until it drops back to $k$. Now every subarray that
ends at $r$ and starts anywhere in $[l,r]$ has at most $k$ distinct values,
contributing $r-l+1$ subarrays. Summing that over all $r$ gives the total.

\begin{ideabox}[Window Width Is the Count of Valid Subarrays]
Because the ``at most $k$ distinct'' property is preserved under shrinking,
once the window $[l,r]$ is valid, so is every subarray ending at $r$ with a
start $\ge l$. Adding $r-l+1$ per right end counts them all without
enumeration.
\end{ideabox}

\subsection*{Algorithm}

\begin{enumerate}
  \item $l=0$, empty frequency map, $\textit{distinct}=0$,
        $\textit{ans}=0$.
  \item For $r=0\ldots n-1$: add $a[r]$; while $\textit{distinct}>k$
        remove $a[l]$ and advance $l$; add $r-l+1$ to $\textit{ans}$.
  \item Output $\textit{ans}$.
\end{enumerate}

\subsection*{Dry run}

\dryrun{$a=[1,2,1,3]$, $k=2$.}

\begin{itemize}
  \item $r=0$: window $[1]$, add $1$. $r=1$: $[1,2]$, add $2$. $r=2$:
        $[1,2,1]$ (distinct $2$), add $3$.
  \item $r=3$: adding $3$ gives distinct $3$; shrink to $[2,1,3]$? still
        $3$; shrink to $[1,3]$, distinct $2$; add $2$.
  \item Total $=1+2+3+2=8$ subarrays with at most $2$ distinct.
\end{itemize}

\subsection*{C++ implementation}

\begin{lstlisting}
#include <bits/stdc++.h>
using namespace std;

int main() {
    int n, k;
    cin >> n >> k;
    vector<int> a(n);
    for (int& x : a) cin >> x;

    unordered_map<int,int> freq;
    long long ans = 0;
    int l = 0, distinct = 0;
    for (int r = 0; r < n; r++) {
        if (freq[a[r]]++ == 0) distinct++;
        while (distinct > k) {
            if (--freq[a[l]] == 0) distinct--;
            l++;
        }
        ans += r - l + 1;                      // subarrays ending at r
    }
    cout << ans << "\n";
    return 0;
}
\end{lstlisting}

\begin{complexitybox}
\textbf{Time Complexity.} $O(n)$ amortised -- each pointer moves forward at
most $n$ times. \textbf{Space Complexity.} $O(k)$ distinct keys.
\end{complexitybox}

\begin{takeawaybox}
``Count subarrays with at most $k$ (property)'' is a variable window where
each right end contributes its window width. To count \emph{exactly} $k$
distinct, take atMost$(k)-$atMost$(k-1)$ -- a standard follow-up on the same
routine.
\end{takeawaybox}

\section{Playlist}
\label{sec:cses-playlist}

\problemstatement{Given a playlist of $n$ songs (each an integer id), find
the length of the longest contiguous stretch of songs with \emph{no
repeats} (all values distinct).}

\begin{patternbox}
``Longest window with all distinct'' is a \textbf{variable two-pointer
window}: extend the right end, and whenever a duplicate appears, jump the
left end past the duplicate's previous occurrence so the window is distinct
again. Track the best width seen.
\end{patternbox}

\subsection*{Jump the left end past duplicates}

Maintain a map from each value to its most recent index. Advance $r$; if
$a[r]$ was seen at some index $\ge l$, move $l$ to one past that index --
this discards exactly the earlier copy and everything before it, restoring
distinctness in one jump rather than shrinking one step at a time. After
each step the window $[l,r]$ is duplicate-free, so update the best length
with $r-l+1$.

\begin{ideabox}[Last-Seen Index Enables an O(1) Left Jump]
Instead of shrinking the window element by element, remembering each value's
last position lets the left pointer leap directly to the smallest valid
start. Both pointers still only move forward, keeping the scan linear.
\end{ideabox}

\subsection*{Algorithm}

\begin{enumerate}
  \item $l=0$, empty last-seen map, $\textit{best}=0$.
  \item For $r=0\ldots n-1$: if $a[r]$'s last index $\ge l$, set
        $l=\text{lastIndex}+1$; update $\text{last}[a[r]]=r$; update
        $\textit{best}=\max(\textit{best},r-l+1)$.
  \item Output $\textit{best}$.
\end{enumerate}

\subsection*{Dry run}

\dryrun{$a=[1,2,1,3,2]$.}

\begin{itemize}
  \item $r=2$ ($a=1$): previous $1$ at index $0\ge l=0$, jump $l=1$;
        window $[2,1]$.
  \item $r=4$ ($a=2$): previous $2$ at index $1<l$? $l$ is now $1$, so
        $1\ge1$, jump $l=2$; window $[1,3,2]$, length $3$.
  \item Best length $=3$.
\end{itemize}

\subsection*{C++ implementation}

\begin{lstlisting}
#include <bits/stdc++.h>
using namespace std;

int main() {
    int n;
    cin >> n;
    vector<int> a(n);
    for (int& x : a) cin >> x;

    unordered_map<int,int> last;               // value -> last index seen
    int l = 0, best = 0;
    for (int r = 0; r < n; r++) {
        auto it = last.find(a[r]);
        if (it != last.end() && it->second >= l) l = it->second + 1;  // jump past dup
        last[a[r]] = r;
        best = max(best, r - l + 1);
    }
    cout << best << "\n";
    return 0;
}
\end{lstlisting}

\begin{complexitybox}
\textbf{Time Complexity.} $O(n)$ amortised. \textbf{Space Complexity.}
$O(n)$ for the last-seen map.
\end{complexitybox}

\begin{takeawaybox}
Longest-distinct-window is a two-pointer scan where a last-seen map lets the
left end jump directly past a repeat. The same ``remember last position,
jump the left bound'' idiom solves longest-substring-without-repeat and many
uniqueness-constrained window problems.
\end{takeawaybox}

\section{Maximum Subarray Sum}
\label{sec:cses-maximum-subarray-sum}

\problemstatement{Given an array of $n$ integers (possibly negative), find
the maximum sum of any non-empty contiguous subarray.}

\begin{patternbox}
The maximum-subarray problem is solved by \textbf{Kadane's algorithm}: a
single scan maintaining the best subarray sum \emph{ending at the current
position}, resetting whenever carrying the previous sum forward would only
hurt.
\end{patternbox}

\subsection*{Best sum ending here}

Let $\textit{cur}$ be the maximum sum of a subarray ending exactly at the
current index. Extending to the next element, either continue the previous
subarray ($\textit{cur}+a[i]$) or start fresh at $a[i]$ -- whichever is
larger:
\[
  \textit{cur}=\max(a[i],\ \textit{cur}+a[i]).
\]
The overall answer is the maximum $\textit{cur}$ over all positions.
Starting fresh happens exactly when the running sum has gone negative and
would drag down any extension.

\begin{ideabox}[Drop a Negative Prefix]
A prefix with negative sum can never help a later subarray -- discarding it
(resetting $\textit{cur}$ to $a[i]$) is always at least as good. This single
greedy reset is what makes one linear pass optimal.
\end{ideabox}

\subsection*{Algorithm}

\begin{enumerate}
  \item $\textit{cur}=\textit{best}=a[0]$.
  \item For $i=1\ldots n-1$: $\textit{cur}=\max(a[i],\textit{cur}+a[i])$;
        $\textit{best}=\max(\textit{best},\textit{cur})$.
  \item Output $\textit{best}$.
\end{enumerate}

\subsection*{Dry run}

\dryrun{$a=[-1,3,-2,5,-4]$.}

\begin{itemize}
  \item $\textit{cur}$: $-1\to3\to1\to6\to2$. Resets never beat carrying
        here except at the start.
  \item $\textit{best}=6$ (the subarray $[3,-2,5]$).
\end{itemize}

\subsection*{C++ implementation}

\begin{lstlisting}
#include <bits/stdc++.h>
using namespace std;

int main() {
    int n;
    cin >> n;
    vector<long long> a(n);
    for (auto& x : a) cin >> x;

    long long cur = a[0], best = a[0];
    for (int i = 1; i < n; i++) {
        cur = max(a[i], cur + a[i]);           // extend or restart
        best = max(best, cur);
    }
    cout << best << "\n";
    return 0;
}
\end{lstlisting}

\begin{complexitybox}
\textbf{Time Complexity.} $O(n)$ -- a single pass. \textbf{Space
Complexity.} $O(1)$.
\end{complexitybox}

\begin{takeawaybox}
Kadane's algorithm is the definitive maximum-subarray solution: keep the
best sum ending here, restart whenever the running total turns negative. Its
``discard a harmful prefix'' idea recurs in maximum-product subarray and
several DP-on-a-line problems.
\end{takeawaybox}

